\scanpage{34}
entièrement explicite et satisfaisante, puisque au deuxième membre interviennent les quantités $u(g'_0),u(g'_1),u(g'_\infty)$, qui ne sont donc pas exprimées directement en termes des \inserted{seules} données $(p,\sigma,g_0,g_1,g_\infty)$ $(p',\sigma',g'_0,g'_1,g'_\infty)$; \deleted{sauf}\inserted{mais ils s'explicitent aussi} quand $g'_0,g'_1,g'_\infty$ sont eux-mêmes donnés explicitement comme des mots (ou limites de tels) en les $l_0,l_1,l_\infty$.

Cependant quand $\tilde u$ est de la forme $i(g,\alpha_0,\alpha_1,\alpha_\infty)$, on trouve (en laissant tomber les $'$ du 2e facteur)
\[
(9)\quad i(g,\alpha_0,\alpha_1,\alpha_\infty)(p,\sigma,g_0,g_1,g_\infty)=(p,\sigma,\deleted{\unsure{\operatorname{int}}}gg_0l_{\unsure{\sigma0|\sigma(0)}}^{\alpha_{\sigma(0)}},gg_1l_{\sigma1}^{\alpha_{\sigma(1)}},gg_\infty l_{\sigma\infty}^{\alpha_{\unsure{\sigma(\infty)}}}\unsure{)}
\]
qui se simplifie pour $\sigma=1$ en
\[
(9')\quad i(g,\alpha_0,\alpha_1,\alpha_\infty)(p,1,g_0,g_1,g_\infty)=(p,\sigma,gg_0l_0^{\alpha_0},gg_1l_1^{\alpha_1},gg_\infty l_\infty^{\alpha_\infty})
\]
Par ailleurs, l'opération de $\widehat{\widehat{\mathcal E}}$ sur $i(\pi\times\mathbb Z^{\{0,1,\infty\}})$ par autom. int. s'explicite aisément, et ne dépend que de $u$ \deleted{\unsure{et}}:
\[
(10)\quad\tilde u\,i(g;\alpha_0,\alpha_1,\alpha_\infty)\tilde u^{-1}=i(u(g);p\alpha_{\sigma^{-1}0},p\alpha_{\sigma^{-1}1},p\alpha_{\sigma^{-1}\infty})
\]
où $\tilde u=(p,\sigma,g_0,g_1,g_\infty)$

--- mais sur cette formule, on peut faire les mêmes commentaires que plus haut, sur le manque d'explicitation du \unsure{2ième|2ème} mb de cette formule, pour $g$ non exprimé comme un mot p.r. aux $l_0,l_1,l_\infty$ (ou limite de tels mots)
\mspage{33}
